% Gradient field of a scalar potential.
% Contour lines of Phi(x,y) = x^2 + 2 y^2 with arrows for grad Phi.
\begin{tikzpicture}[>=Stealth, scale=1.0]
  \begin{axis}[
    width=10cm, height=7cm,
    xlabel={$x$}, ylabel={$y$},
    xmin=-2.2, xmax=2.2, ymin=-1.6, ymax=1.6,
    axis lines=center,
    label style={font=\footnotesize},
    tick label style={font=\footnotesize},
    samples=21,
  ]
    % Contours of Phi = x^2 + 2 y^2.
    \addplot[engblue, very thin, domain=0:360, samples=80] ({0.6*cos(x)}, {0.6/sqrt(2)*sin(x)});
    \addplot[engblue, very thin, domain=0:360, samples=80] ({1.0*cos(x)}, {1.0/sqrt(2)*sin(x)});
    \addplot[engblue, very thin, domain=0:360, samples=80] ({1.4*cos(x)}, {1.4/sqrt(2)*sin(x)});
    \addplot[engblue, very thin, domain=0:360, samples=80] ({1.8*cos(x)}, {1.8/sqrt(2)*sin(x)});

    % Gradient arrows: grad Phi = (2x, 4y); arrows scaled.
    \addplot[
      black,
      quiver={u=2*x/12, v=4*y/12, scale arrows=1.0, every arrow/.append style={-Stealth, thin}},
    ] coordinates {
      (-1.5,-1) (-1,-1) (-0.5,-1) (0.0,-1) (0.5,-1) (1,-1) (1.5,-1)
      (-1.5,-0.5) (-1,-0.5) (-0.5,-0.5) (0.0,-0.5) (0.5,-0.5) (1,-0.5) (1.5,-0.5)
      (-1.5,0.5) (-1,0.5) (-0.5,0.5) (0.0,0.5) (0.5,0.5) (1,0.5) (1.5,0.5)
      (-1.5,1) (-1,1) (-0.5,1) (0.0,1) (0.5,1) (1,1) (1.5,1)
    };

    \node[font=\footnotesize, anchor=west, engblue] at (axis cs:1.5,0.05) {level sets of $\Phi$};
    \node[font=\footnotesize, anchor=west] at (axis cs:0.7,1.25) {$\grad\Phi$ (perpendicular to level set)};
  \end{axis}
\end{tikzpicture}
